\section{弦在图像上方意味着什么}
\label{sec:09-Convex-Optimization/03-Convex-Functions/01}

一个回归模型上线一周，监控面板上的平均平方误差稳定在 \(0.84\)。有人为了省一次遍历改了统计口径：先把这一批样本的残差平均起来，再平方。指标掉到 \(0.11\)。模型没动，数据没动，动的只是两个运算的先后。

设第 \(i\) 条样本的残差是 \(r_i\)。原口径算的是 \(\frac1n\sum_i r_i^2\)，新口径算的是 \(\big(\frac1n\sum_i r_i\big)^2\)，两者之差

\[
\frac1n\sum_i r_i^2-\Big(\frac1n\sum_i r_i\Big)^2
\]

正是这批残差的样本方差，非负。只要残差不是完全相同，新口径就严格偏小。指标好看了，误差一点没少。

错的不在代码。平方与平均之间存在一个固定方向的偏差，它跟残差取什么值无关，只跟平方函数的形状有关。把这个形状单独拎出来看。

在 \(f(x)=x^2\) 的图像上取 \((1,1)\) 与 \((3,9)\)，连一条弦。\(x=2\) 处弦的高度是 \(0.5\cdot1+0.5\cdot9=5\)，而 \(f(2)=4\)。换成 \((0,0)\) 与 \((4,16)\)，\(x=2\) 处弦高 \(8\)，函数值还是 \(4\)。换不等的权重也一样：\(\theta=0.3\) 时弦高 \(0.7\cdot1+0.3\cdot9=3.4\)，而 \(f(1.6)=2.56\)。弦始终压在图像上方。

\begin{center}
\begin{tikzpicture}[x=1.35cm,y=0.52cm,>=stealth]
  \draw[->] (-0.25,0) -- (3.65,0) node[right] {\(x\)};
  \draw[->] (0,-0.4) -- (0,10.4) node[above] {\(f(x)\)};
  \draw[blue!70!black,thick] plot[domain=0:3.25,samples=60] (\x,{\x*\x});
  \draw[black,thick] (1,1) -- (3,9);
  \fill (1,1) circle (2.2pt);
  \fill (3,9) circle (2.2pt);
  \fill[blue!70!black] (2,4) circle (2.2pt) node[below right] {\(f(2)=4\)};
  \fill (2,5) circle (2.2pt) node[above left] {弦高 \(=5\)};
  \draw[<->,gray!80] (2.12,4) -- (2.12,5);
\end{tikzpicture}

\emph{凸性比较的是两种顺序：先混合输入再计算函数，落在曲线上；先计算端点再混合函数值，落在弦上。}
\end{center}

监控面板上那两个数字，就是图里竖直相隔的那一小段。弦上的点是先平方再平均，曲线上的点是先平均再平方，\(0.84\) 与 \(0.11\) 分别站在这一段的两头，中间那截长度等于残差的方差。换任何一个开口向上的损失函数，方向都不会反过来。

方向会反过来的函数当然存在。取 \(f(x)=\sin x\)，端点选 \((\pi/4,\sqrt2/2)\) 与 \((3\pi/4,\sqrt2/2)\)：中点 \(x=\pi/2\) 处弦高 \(0.707\)，而 \(\sin(\pi/2)=1\)，弦跑到了图像下方。可换成 \((\pi/2,1)\) 与 \((3\pi/2,-1)\)，中点处弦高 \(0\)，\(\sin\pi=0\)，两者持平。同一个 \(\sin x\)，有的弦在上，有的弦在下，有的相切。它没有固定的偏差方向，因而不是凸函数。

\subsection{弦不等式就是凸函数的定义}

上面反复做的那件事，写成一行就是定义。设 \(f:\R^n\to\R\) 的定义域是凸集。若对定义域中任意 \(x,y\) 和任意 \(0\le\theta\le1\)，

\[
f(\theta x+(1-\theta)y)
\le
\theta f(x)+(1-\theta)f(y),
\]

则称 \(f\) 为\textbf{凸函数}。

这里没有极限、没有导数、没有连续性假设，只有加权平均和一个不等号。左边是先在输入空间里混合、再评价；右边是先评价两个端点、再混合评价结果。凸性给出的断言是：混合方案的代价不会超过端点代价的线性插值。用上一单元投资组合的语言说，两种配置各有成本，按比例混合后的成本不会高过两种成本的加权平均——混合不额外收费。可行域那一半的保证在\hyperref[sec:09-Convex-Optimization/02-Convex-Sets/01]{``线段、凸组合与凸包''}里已经建立，目标函数这一半从这里开始。

\subsection{定义域必须和公式一起给出}

判断凸性不能只盯着表达式。\(1/x\) 在正半轴 \((0,\infty)\) 上二阶导数为正，是凸的；在负半轴 \((-\infty,0)\) 上二阶导数为负，是凹的（二阶导数判据在\hyperref[sec:09-Convex-Optimization/03-Convex-Functions/03]{``曲率区分凸、严格凸与强凸''}一节严格建立）。若把定义域写成 \(\R\setminus\{0\}\)，麻烦出在更前面一层：取 \(x=-1\) 与 \(y=1\)，凸组合 \(0\) 不在定义域内，弦不等式的左边根本没有定义。定义域本身不凸，函数凸不凸就无从谈起。

优化里常用\textbf{扩展值函数}把约束吸收进目标：

\[
I_C(x)=
\begin{cases}
0,&x\in C,\\
+\infty,&x\notin C.
\end{cases}
\]

当 \(C\) 是凸集时，指标函数 \(I_C\) 在扩展值意义下是凸函数。于是带约束的问题 \(\min_{x\in C}f(x)\) 可以等价写成

\[
\min_x\ f(x)+I_C(x),
\]

形式上没有约束了，可行性信息一点没丢，只是搬进了目标函数。这不是记号游戏：近端方法处理的正是``光滑项加一个可能取 \(+\infty\) 的凸项''这种结构，见\hyperref[sec:09-Convex-Optimization/06-Optimization-Algorithms/04]{``近端与坐标方法利用可分结构''}。

\subsection{Jensen 不等式把两点推广到任意平均}

弦不等式只管两个点。对点数做归纳，立刻推广到任意有限凸组合：

\[
f\left(\sum_{i=1}^k\theta_i x_i\right)
\le
\sum_{i=1}^k\theta_i f(x_i),
\qquad
\theta_i\ge0,\quad\sum_i\theta_i=1.
\]

权重是一个概率分布，于是概率形式顺理成章：只要期望存在且 \(X\) 落在定义域内，

\[
f(\E[X])\le\E[f(X)].
\]

这就是 Jensen 不等式，它的严格版本、取等条件和随机情形下的细节见\hyperref[sec:05-Probability/07-Transforms-and-Bounds/03]{``Jensen 不等式与随机凸性''}。开头那个监控指标是它的 \(k=n\)、等权重、\(f(x)=x^2\) 的特例。

同一个不等式换几个 \(f\) 就换一副面孔。取 \(f(x)=x^2\) 得 \((\E X)^2\le\E[X^2]\)，移项就是 \(\Var(X)=\E[X^2]-(\E X)^2\ge0\)——方差非负是 Jensen 的一句注脚。取 \(f(x)=-\log x\)（\(x>0\) 上凸），代进去整理即得加权算术--几何平均不等式。取 \(f(x)=e^x\)，得 \(\exp(\E[X])\le\E[e^X]\)，这是 Chernoff 型尾概率界的起点。

方向朝下的用法更值钱。\(\log\) 是凹函数，Jensen 反向：\(\log\E[Z]\ge\E[\log Z]\)。变分推断和 VAE 里的证据下界就是这么来的——边缘似然 \(\log p(x)=\log\E_{q}[p(x,z)/q(z)]\) 里那个积分算不动，把 \(\log\) 换到期望里面，得到一个算得动的下界，然后去最大化它。凹性保证了换过去只会变小，因此最大化下界是一件安全的事。这条路线的完整推导在\hyperref[sec:13-Machine-Learning/08-Sampling-and-Inference/03]{``变分推断把后验近似变成优化''}与\hyperref[sec:14-Deep-Learning/06-Variational-Autoencoders/01]{``VAE 把不可积后验变成可训练下界''}。

每次看到一个期望在非线性函数的里外挪位置，Jensen 都在告诉你不等号朝哪边开口。

\subsection{上境图把函数问题换成集合问题}

函数不好直接摆弄的时候，可以改看它的上方区域。定义 \(f\) 的\textbf{上境图}为

\[
\operatorname{epi}f=\set{(x,t):t\ge f(x)}.
\]

它由图像本身和图像上方的所有点组成。关键的等价关系是：\(f\) 是凸函数当且仅当 \(\operatorname{epi}f\) 是凸集。

两个方向都只需要一步。设 \((x,t)\) 与 \((y,s)\) 在上境图中，即 \(t\ge f(x)\)、\(s\ge f(y)\)，凸性给出

\[
f(\theta x+(1-\theta)y)
\le
\theta f(x)+(1-\theta)f(y)
\le
\theta t+(1-\theta)s,
\]

混合点仍在上境图内。反向取边界点 \((x,f(x))\) 与 \((y,f(y))\)，混合点落在上境图内这句话展开就是弦不等式。

换成集合以后，函数的构造规则变成了集合的交、仿射像和投影，而这些操作在上一单元已经熟悉。举一个立刻见效的例子：逐点最大值 \(f(x)=\max_i f_i(x)\) 的上境图是各个 \(\operatorname{epi}f_i\) 的交集，交集保凸，所以只要每个 \(f_i\) 凸，最大值就凸。一行论证覆盖了 hinge 损失、\(\ell_\infty\) 范数和一大类鲁棒目标，下一步的推广见\hyperref[sec:09-Convex-Optimization/03-Convex-Functions/04]{``怎样识别和构造凸函数''}。

\subsection{次水平集只提供单向判据}

凸函数的次水平集 \(\set{x:f(x)\le\alpha}\) 一定是凸集。若 \(x,y\) 都在其中，

\[
f(\theta x+(1-\theta)y)
\le
\theta f(x)+(1-\theta)f(y)
\le
\theta\alpha+(1-\theta)\alpha=\alpha.
\]

反过来不成立。所有次水平集都凸的函数叫\textbf{拟凸函数}，它可以完全不满足弦不等式。一元单调函数就是现成的反例：次水平集全是区间，因而全凸，但函数本身可以任意弯曲。

这个区别在读等高线图时容易被跳过。等高线看着像一圈圈的碗，只说明``不太差''的那部分区域是连通且不带凹口的，不保证任何一条弦不会穿进图像里去。拟凸问题有自己的算法（二分加可行性判定），但它拿不到下一篇要建立的那条下界。

\subsection{几块积木先记住}

以下函数在各自的自然定义域上凸，后面构造复杂目标时会反复用到：

\begin{itemize}
\tightlist
\item
  仿射函数 \(a^{\trans}x+b\)，弦不等式恒取等号，它既凸又凹；
\item
  任意范数 \(\norm{x}\)，凸性由三角不等式与齐次性直接给出；
\item
  二次函数 \(\frac12x^{\trans}Qx+a^{\trans}x+b\)，当且仅当 \(Q\succeq0\) 时凸；
\item
  指数 \(e^x\) 在整条实轴上严格凸，负对数 \(-\log x\) 在 \(x>0\) 上严格凸，后者是最大似然与信息论里最常见的那一个；
\item
  log-sum-exp \(\log\sum_i e^{x_i}\)，光滑、严格凸，是最大值函数的软化版本，softmax 的分母就是它。
\end{itemize}

反面同样要记：\(-x^2\) 是凹的，弦落在图像下方；\(\sin x\) 在整条实轴上既不凸也不凹，只在某些区间上有单一的曲率方向。判断凸性时把定义域一并写出，这个习惯能挡掉后面很多莫名其妙的报错。

到这里，凸性还是一条关于两个点的陈述：要用它，得同时拿着 \(x\) 和 \(y\)。但优化过程中我们站在一个点上，手里只有这一点的函数值和梯度，另一个点恰恰是要找的东西。下一篇把弦不等式的两端拉到无穷近，看看一个点的局部信息能对整个函数说出什么。
